% Generated draft: requires mathematical and source review.
\documentclass[11pt]{article}
\usepackage[margin=25mm]{geometry}
\usepackage{fontspec}
\setmainfont{Latin Modern Roman}
\usepackage{amsmath,amssymb,mathtools,unicode-math}
\setmathfont{Latin Modern Math}
\usepackage{enumitem,needspace,xurl,hyperref}
\setlength{\parindent}{0pt}
\setlength{\parskip}{5pt}
\setlength{\emergencystretch}{3em}
\newcommand{\N}{\mathbb{N}}
\newcommand{\Z}{\mathbb{Z}}
\newcommand{\Q}{\mathbb{Q}}
\newcommand{\R}{\mathbb{R}}
\newcommand{\C}{\mathbb{C}}
\newcommand{\F}{\mathbb{F}}
\newcommand{\A}{\mathbb{A}}
\newcommand{\PP}{\mathbb{P}}
\newcommand{\E}{\mathbb{E}}
\newcommand{\eps}{\varepsilon}
\newcommand{\dd}{\,\mathrm{d}}
\begin{document}
\section*{UnsolvedMath ID: 30001758}
\textbf{Source problem code:} catalogue entry
\subsection{Definitions and mathematical statement}
A closed manifold is a compact manifold without boundary. Write \(T^k=(S^1)^k\) for the \(k\)-dimensional torus, where \(k\in\Z_{\geq 0}\) and \(T^0\) is a point. Define the first rational Betti number by
\[
b_1(X)=\dim_{\Q}H^1(X;\Q).
\]
For a nonempty closed manifold \(X\), its rational cuplength is
\[
\operatorname{cup}_{\Q}(X)=\max\Bigl(\{0\}\cup\bigl\{r\in\Z_{\geq 1}:\text{there exist }d_1,\ldots,d_r\geq 1\text{ and }\alpha_i\in H^{d_i}(X;\Q)\text{ with }\alpha_1\smile\cdots\smile\alpha_r\neq 0\bigr\}\Bigr).
\]
The product-splitting property requested in the supplied record is the existence of a finite-sheeted covering \(p:\widetilde M\to M\), an integer \(k\geq 0\), and a simply connected manifold \(N\) such that
\[
\widetilde M\cong T^k\times N,
\qquad
b_1(M)\leq k\leq\operatorname{cup}_{\Q}(\widetilde M).
\]
Here \(\cong\) is a placeholder for the intended type of product splitting: the supplied extraction does not specify whether it means a homeomorphism, a diffeomorphism, or an isometry with a product metric. Simply connected means path connected with trivial fundamental group.

\textbf{Question.} Which curvature hypotheses on closed manifolds ensure that every manifold satisfying those hypotheses has the displayed finite-cover product-splitting property? The original formulation additionally asks for hypotheses beyond the known cases and for further topological obstructions to the requested conclusion.

\textbf{Unresolved qualifications.} The supplied record does not identify the known cases, specify the curvature hypotheses or metric quantifiers, state connectedness and regularity assumptions, or fix the category of splitting. Curvature suggests a Riemannian setting, but that interpretation and its precise scope require confirmation. The intended target of the additional topological-obstruction question is also unspecified. Consequently, the displayed conclusion records the explicit target supplied by the extraction; it does not supply a complete conjecture with fixed hypotheses.
\subsection{Short English statement}
Identify curvature hypotheses forcing a closed manifold to have a finite cover that is a product of a torus and a simply connected manifold, with the torus dimension bounded below by \(b_1(M)\) and above by the cover's rational cuplength. The source also asks for additional topological obstructions. The known cases, precise curvature assumptions, metric quantifiers, and meaning of splitting remain unspecified.
\subsection{Sources}
\begin{itemize}
\item Supplied local UnsolvedMath registry record, permanent ID \(30001758\), title: \emph{Curvature-Driven Product Splittings and Topological Obstructions}; exact source problem code: \texttt{catalogue entry}. Collection: \emph{Oberwolfach Reports Open Problems}. The record gives the year 2011 and names no individual proposer.
\item Supplied report citation: \emph{Arbeitsgemeinschaft: Rational Homotopy Theory in Mathematics and Physics} (2011), Oberwolfach Reports, DOI 10.4171/owr/2011/18, pp. 993--1045. \url{https://doi.org/10.4171/owr/2011/18}. The background attributes the problem to workshop open-problem discussions documented in Oberwolfach Reports; no individual problem page is supplied.
\item Supplied literature source: \emph{Rational Homotopy Theory in Mathematics and Physics, OWR 18/2011}. \url{https://doi.org/10.4171/owr/2011/18}.
\item Literature assessment checked at \texttt{2026-08-21}: the supplied curation labels the problem open and reports that the extraction supplies no precise theorem and that no general resolution was identified. This is supplied dated curation, not an independently verified literature assessment. The background expressly requests independent verification of time-sensitive status and cited claims.
\item The registry labels the statement \texttt{corrected\_verified} and reports that its mathematical content was checked against the official report, with notation, context, and question form normalized without changing the source-level claim. This verification claim is supplied registry metadata and has not been independently checked here.
\end{itemize}
\end{document}
